\section{Multi-GPU Computing with Julia} 

Many Julia applications can use multiple GPUs simultaneously. Typical examples include: 

\begin{itemize} 
\item Oceananigans.jl 
\item MPI-based PDE solvers 
\item Distributed machine learning 
\item Large ensemble simulations 
\end{itemize} 

Multi-GPU programs typically require a CUDA-aware MPI installation. The typical workflow is: 

\begin{enumerate} 
\item Load CUDA, OpenMPI, and Julia modules. 
\item Configure MPI.jl to use the system MPI installation. 
\item Launch Julia using \texttt{mpiexec}. 
\item Verify that each MPI rank can access a GPU. \end{enumerate} 

\subsection{Loading Required Modules} 

Load the required software: 

\begin{lstlisting} 
module load cuda/12.3.2 
module load openmpi/4.1.7-with-cuda 
module load julia/1.10.8 
\end{lstlisting} 

Verify that the modules are loaded: 

\begin{lstlisting} 
module list 
\end{lstlisting} 

Verify CUDA: 

\begin{lstlisting} 
nvidia-smi 
\end{lstlisting} 

Verify MPI: 

\begin{lstlisting} 
which mpiexec 
\end{lstlisting}


\subsection{Installing Required Julia Packages} 

The examples in this appendix use the following Julia packages: 

\begin{itemize} 
\item CUDA.jl 
\item MPI.jl 
\item MPIPreferences.jl 
\end{itemize} 

If they are not already installed, start Julia and run: 


\begin{lstlisting} 
using Pkg 
Pkg.add("CUDA") 
Pkg.add("MPI") 
Pkg.add("MPIPreferences") 
Pkg.add("Oceananigans") 
\end{lstlisting}

\subsection{Testing CUDA} 

Before testing MPI, verify that Julia can see a CUDA-enabled GPU. Start Julia: 

\begin{lstlisting} 
julia 
\end{lstlisting} 

Load CUDA.jl and check for GPUs: 

\begin{lstlisting} 
using CUDA 
CUDA.has_cuda_gpu() 
\end{lstlisting} 

Expected output: 

\begin{lstlisting} 
true 
\end{lstlisting} 

You may also display information about the available GPUs: 

\begin{lstlisting} 
CUDA.devices() 
\end{lstlisting}

\subsection{Configuring MPI.jl} 

Inside Julia: 

\begin{lstlisting} 
using MPIPreferences 
MPIPreferences.use_system_binary() 
\end{lstlisting} 

Restart Julia and verify: 

\begin{lstlisting} 
using MPI 
MPI.versioninfo() 
\end{lstlisting}

\subsection{Verifying GPU Assignment} 

Create a file called \texttt{test.jl}: 

\begin{lstlisting} 
using MPI 
using CUDA 
MPI.Init() 
rank = MPI.Comm_rank(MPI.COMM_WORLD) 
println( "Rank $rank using GPU $(CUDA.device())" ) 
MPI.Finalize() 
\end{lstlisting} 

Launch four MPI ranks: 

\begin{lstlisting} 
mpiexec -n 4 julia test.jl 
\end{lstlisting}

Expected output: 

\begin{lstlisting} 
Rank 0 using GPU CuDevice(0) 
Rank 1 using GPU CuDevice(1) 
Rank 2 using GPU CuDevice(2) 
Rank 3 using GPU CuDevice(3) 
\end{lstlisting}


\subsection{Selecting GPUs} 

On a shared server, it is often desirable to restrict a program to a subset of the available GPUs. Before launching a computation, check current GPU usage: 

\begin{lstlisting} 
nvidia-smi 
\end{lstlisting} 

You can restrict a program to specific GPUs using \texttt{CUDA\_VISIBLE\_DEVICES}. Use only GPU 0: 

\begin{lstlisting} 
CUDA_VISIBLE_DEVICES=0 julia script.jl 
\end{lstlisting} 

Use GPUs 1 and 2: 

\begin{lstlisting} 
CUDA_VISIBLE_DEVICES=1,2 julia script.jl 
\end{lstlisting} 

The selected GPUs are re-numbered internally. For example, 

\begin{lstlisting} 
CUDA_VISIBLE_DEVICES=2 julia script.jl 
\end{lstlisting} 

makes physical GPU 2 appear as GPU 0 inside Julia. This is useful on shared servers where multiple users may be using GPUs simultaneously. When launching MPI jobs, the same variable can be used: 

\begin{lstlisting} 
CUDA_VISIBLE_DEVICES=0,1,2,3 mpiexec -n 4 julia test.jl 
\end{lstlisting} 

Each MPI rank can then be assigned a GPU programmatically.

\subsection{Checking Available GPUs} 

Before launching a large computation, check current GPU usage: 

\begin{lstlisting} 
nvidia-smi 
\end{lstlisting} 

This allows you to identify which GPUs are currently in use.

\subsection{Checking Which GPU Julia Is Using} 

Inside Julia: 

\begin{lstlisting} 
using CUDA 
CUDA.device() 
\end{lstlisting}